% ============================================================
%  CAPÍTULO 8: LIMITACIONES
% ============================================================
\chapter{Limitaciones del estudio}
\label{cap:limitaciones}

% ------------------------------------------------------------
\section{Limitaciones de los datos}
% ------------------------------------------------------------

El análisis empírico del presente trabajo está sujeto a un conjunto de
limitaciones que deben tenerse en cuenta al interpretar los resultados.

\subsection{Disponibilidad y calidad estadística de las fuentes}

La Encuesta de Presupuestos Familiares (EPF) utilizada como fuente de la variable
dependiente mide el gasto \textbf{declarado} por los hogares, lo que puede
subestimar el gasto real si los encuestados olvidan o infravaloran el coste de las
primas de seguro de salud incluidas en nómina o en coberturas empresariales. La
EPF tampoco distingue entre seguros de salud individuales y colectivos, lo que
impide aislar el componente de decisión voluntaria del hogar.

\subsection{Período temporal y variación}

El panel cubre 2016--2024, un período de nueve años que incluye cambios
significativos (pandemia 2020, reforma de listas de espera 2022). La limitación del
horizonte temporal restringe el número de observaciones disponibles para las
estimaciones con retardos ($T-2 = 7$ períodos efectivos) y limita la potencia
estadística de los contrastes de raíz unitaria y cointegración. Una extensión del
panel hasta 2030, actualmente no disponible, permitiría mayor precisión en la
estimación del multiplicador de largo plazo.

\subsection{Ausencia de datos de precio de las primas}

El modelo teórico predice que el precio relativo del seguro privado respecto al
coste del servicio público es el mecanismo central de decisión. Sin embargo, el
precio de las primas de seguros sanitarios no está disponible a nivel de CCAA en
series temporales homologables. La omisión de esta variable puede generar un sesgo
de variable omitida si el precio de los seguros está correlacionado con el PIB
per cápita o con los tiempos de espera.

\subsection{Heterogeneidad en los sistemas de espera}

Los datos de SISLE-SNS presentan heterogeneidad en la metodología de registro
entre CCAA. Algunas comunidades registran las esperas de forma continua mientras
otras lo hacen puntualmente, lo que puede introducir ruido de medición en la
variable $\text{Espera}_{it}$. Este error de medición, si es clásico, sesga los
coeficientes hacia cero (atenuación), lo que implica que el efecto real de la
espera sobre el gasto en seguros sería \textbf{mayor} que el estimado.

% ------------------------------------------------------------
\section{Limitaciones metodológicas}
% ------------------------------------------------------------

\subsection{Endogeneidad potencial}

Aunque los efectos fijos eliminan la heterogeneidad no observada constante en el
tiempo, existe la posibilidad de \textbf{endogeneidad simultánea}: una CCAA con
mayor demanda de seguros privados (debido a características no observadas de su
población) podría ejercer menos presión política para reducir las listas de espera.
Este canal inverso ---seguros $\rightarrow$ espera--- no ha podido instrumentarse
por falta de instrumentos exógenos válidos a nivel de CCAA.

\subsection{Umbral endógeno}

El umbral de 100 días utilizado en M3 es establecido \textit{a priori} con base en
los estándares del SISLE-SNS. Un análisis formal de umbral endógeno
\parencite{Maddala1999} con estimación del punto de ruptura óptimo mejoraría la
especificación del modelo M3, aunque requeriría un número de observaciones mayor
que el disponible para garantizar la potencia estadística.

\subsection{Ausencia de nivel microeconómico}

El análisis trabaja a nivel de CCAA, lo que impide identificar el comportamiento
del hogar individual. Idealmente, se requeriría datos de panel a nivel de hogar
(combinando EPF longitudinal con datos de salud) para contrastar si los hogares
de mayor renta reaccionan más intensamente a los incrementos de listas de espera.

% ------------------------------------------------------------
\section{Implicaciones para futuras investigaciones}
% ------------------------------------------------------------

Las limitaciones identificadas sugieren cuatro líneas de extensión:

\begin{enumerate}
  \item \textbf{Datos de panel de hogares:} Combinar la EPF con el módulo de salud
    de la Encuesta Nacional de Salud (ENS) y datos de SISLE-SNS a nivel de hogar.
  \item \textbf{Instrumentos para la endogeneidad:} Utilizar cambios en la
    financiación autonómica o reformas de gestión hospitalaria como instrumentos
    para las listas de espera.
  \item \textbf{Análisis de umbral endógeno:} Cuando el panel se amplíe a 15
    o más períodos, estimar el punto de ruptura óptimo $\tau^*$ mediante los
    estimadores de \textcite{Maddala1999}.
  \item \textbf{Extensión internacional:} Comparar los resultados con países
    con sistemas tipo NHS similares al español (Italia, Portugal, Grecia) para
    identificar patrones comunes.
\end{enumerate}
